\section{Method: generator-first stability classification}

\subsection{The unified object}

Every supported input reduces to the spectrum of one generator:
\begin{itemize}
\item Second-order (vibrations, phonons, adsorbate modes): \(M\ddot{q} + C\dot{q} + Kq = 0\), solved as the quadratic eigenvalue problem \citep{tisseur}, \((\lambda^2 M + \lambda C + K)u = 0\), via companion linearization, with the left and right eigenvectors of the response layer following the non-self-adjoint spectral construction of \citet{keldysh}.
\item First-order (microkinetics, master equations): \(\dot{z} = Gz\).
\end{itemize}
For a scalar second-order mode the poles are \(\lambda = -\gamma/2 \pm i\sqrt{\omega_0^2 - \gamma^2/4}\), and the damping ratio \(\chi = \gamma/2\omega_0\) separates the complex-pole (oscillatory, \(\chi<1\)) regime from the real-pole (monotone, \(\chi>1\)) regime, with a repeated root at \(\chi=1\).

\subsection{Classification vocabulary}

ChemSA reports, per mode, four orthogonal properties: stability (stable, marginal, unstable); temporal form (oscillatory, monotone, repeated-root; the term critically damped is reserved for a second-order mechanical reading, where it is licensed); structure (simple, semisimple-at-tolerance, defective-at-tol, higher-order-ep-candidate, near-coalescent-unresolved, unresolved-representation), reported alongside provenance and a scoped interpretation; and mechanical role (well, saddle, soft-mode, node). The bare label \texttt{defective} is not produced for any multiplicity: floating-point evidence, however clean, is tolerance-qualified by construction, and \texttt{defective-at-tol} says so; a candidate of order \(m>2\) that cannot be certified even at that qualification is labelled \texttt{higher-order-ep-candidate} instead (Section~2.3(b)). A mode is not forced into a single label; a system with a saddle and a stable oscillation reports both.

These properties belong to one of three layers, which answer different questions and should not be read as one universal scalar:

\begingroup\small
\noindent\begin{tabular}{@{}p{0.10\linewidth}p{0.42\linewidth}p{0.40\linewidth}@{}}
\toprule
Layer & Object & Question \\
\midrule
A: state & poles, mechanical $\chi$ where licensed, spectral pole-angle $\rho$, temporal and multiplicity class & Where is the mode? \\
B: motion & finite pole shift $(-\Delta\mathrm{Re}\,\lambda, \Delta\mathrm{Im}\,\lambda)$ & What changed between two supplied generators? \\
C: response & left and right eigendata, residues, sensitivity, conditioning & How is it excited, observed, and moved? \\
\bottomrule
\end{tabular}
\endgroup

Layer B is descriptive and path-dependent: it denotes the finite difference between two supplied generator spectra and is retained here as a conceptual comparison, not as a supported software module or a validated general real-system diagnostic. Layer C is predictive and channel-dependent: it maps declared perturbations and drive/readout channels to pole motion and observable response (Section~2.5). Neither layer is another universal scalar alongside $\chi$.

\subsection{The safeguards}

Six safeguards distinguish ChemSA from a naive linewidth-to-\(\chi\) conversion, in two categories: three core scientific safeguards, and three implementation validity gates that make the first three trustworthy in practice.

\textbf{Core scientific safeguard 1: exceptional-point versus semisimple discrimination.} This is the core safeguard, and the one a frequency-and-linewidth reading cannot reproduce. Consider two \(2\times2\) generators with the same repeated eigenvalue \(\lambda = -1\):
\begin{itemize}
\item \(G_1 = \mathrm{diag}(-1,-1) = -I\), a repeated root with a full eigenbasis (every vector is an eigenvector). This is a semisimple degeneracy: the modes are independent and separate smoothly under perturbation. In parameterized Hermitian or normal systems this nondefective case is often called a diabolic point; the more general term semisimple is used here, since a repeated scalar matrix is not necessarily a conical diabolic point in the strict sense.
\item \(G_2 = \left[\begin{smallmatrix}-1 & 1\\ 0 & -1\end{smallmatrix}\right]\), the same repeated root \(-1\) but a single eigenvector (geometric multiplicity 1 less than algebraic multiplicity 2). This is a defective matrix: a genuine exceptional point with square-root perturbation sensitivity.
\end{itemize}
Their eigenvalues are identical; only the eigenvectors differ. ChemSA computes algebraic and geometric multiplicity separately and returns \texttt{semisimple-at-tolerance} for \(G_1\) and \texttt{defective-at-tol} for \(G_2\), demonstrating that it measures eigenvector deficiency, not eigenvalue closeness. Its contribution over a textbook Jordan-form check is extending this to the second-order quadratic eigenvalue problem via companion linearization, with the \(\sqrt{\varepsilon}\) tolerance appropriate to EP2 (algebraic-multiplicity-two) defective-eigenvalue splitting -- an order-\(m\) Jordan block generically splits as \(\varepsilon^{1/m}\) instead, the subject of Section~2.3(b) -- and with physical input contracts that reject non-mechanical generators. The discrimination is verified by randomized regression sweeps over constructed matrices, and separately by one physically parameterized model-class cross-check (Section~2.6); neither alone would suffice. 500 randomized-parameter twin oscillators produced zero false exceptional points, and 500 randomized-parameter exact-critical systems produced zero misses (Supplementary~S3, reproducible from the test suite with a fixed seed). The suite is adversarial: the exact cases that exposed defects during development are retained as permanent regression tests. Because a finite-precision eigensolver cannot supply a structural or symbolic certificate of exact defectiveness, even the exact-entry Jordan example is reported as \texttt{defective-at-tol}; the unqualified \texttt{defective} label is not produced by the numerical classifier.

\textbf{Core scientific safeguard 2: provenance gating.} Defectiveness established from a supplied equation is defectiveness of that representation. A reduced or effective generator may have a genuine exceptional point of its own effective dynamics; what it cannot do, without further analysis, is establish defectiveness of the parent generator it came from, since a repeated root produced by eliminating variables can be an artifact of the elimination. This is an interpretive constraint, not a measurement one: the multiplicities do not change because a flag changed. ChemSA therefore reports three separate fields, structure (the numerical finding, provenance-independent), provenance (full, reduced, or unknown, as declared by the caller), and interpretation (what the structure licenses, provenance-dependent), rather than one overloaded label. An earlier version folded these into a single degeneracy label, so the same defective cluster read differently under full and reduced provenance, which erased the measurement in order to state the caveat. The gate remains a mandatory disclosure and interpretation guard, not an automated truth detector: it cannot catch a mislabelled input, only prevent a silent upgrade from reduced to parent.

\textbf{Core scientific safeguard 3: the constructive scope fence.} The third safeguard is a limit on what the engine may be asked for at all. Local generator architecture classifies a well; it does not determine a rate. Section~3 makes this constructive rather than declarative: two systems with IDENTICAL reactant wells, and therefore identical classification by the engine, differ in activated rate by a factor of $e^{5} \approx 148$ through the global barrier alone. The fence is executable, is run as a validator stage, and classifies both wells through the same public engine the rest of this work uses, so a drift in that engine would break it.

\textbf{Implementation validity gate 1: source gating of numerical inputs.} The provenance gate above governs the representation an equation came from; a second, data-side gate governs whether a number is entitled to appear at all. Every quantity entering a reported calculation is carried as a value paired with its source, and a value with no source is held to be absent: it contributes no number, and any derived quantity that depends on it is returned blocked rather than estimated, inheriting the weakest provenance of its inputs. It is a disclosure and blocking mechanism, not a check that a cited source actually supports its value; that remains a human responsibility. The worked example of Part~III exercises it directly: its one quantitative result, the endpoint width growth, is computed through the gate from the two source-anchored widths, so an illustrative digitized point is structurally unable to contribute to it.

\textbf{Implementation validity gate 2: scalar-\(\chi\) applicability.} A single scalar \(\chi\) is meaningful only when the damping is proportional (\(C\) and \(K\) simultaneously diagonalizable in the mass metric, the Caughey--O'Kelly condition \citep{caughey}). This condition is TESTED by the engine, not accepted on declaration: \texttt{scalar\_reduction\_valid} computes the relative mass-metric commutator of \(M^{-1}C\) and \(M^{-1}K\) and withholds the mechanical \(\chi\), with a stated reason, whenever it is nonzero. When it is not, ChemSA reports not-applicable rather than forcing a number.

\textbf{Implementation validity gate 3: finite-precision equality and certified-domain caution.} A tolerance cluster is not a proof of exact equality, in either direction. Just as a numerically coalescent pair is labelled defective-at-tolerance rather than a certified exceptional point, a repeated eigenvalue with a full eigenbasis is labelled semisimple-at-tolerance, never as exact semisimplicity, which finite precision cannot establish. The reserved label semisimple-exact is used only for multiplicity established analytically or structurally, and the engine never asserts it from a float computation.

\subsection{Scale and representation invariance}

Quadratic eigensolution is performed in mass-normalized coordinates, which removes ordinary coordinate-unit scaling and improves congruence robustness. The certified numerical range is empirical and limited by conditioning: the release reports the mass condition number, applies the tested acceptance gates, and refuses cases outside the calibrated domain. No finite-precision implementation is claimed to be invariant under arbitrarily ill-conditioned congruences; a random non-diagonal congruence was measured to defeat exact-coincidence certification at a mass condition number as low as \(10^2\), well below machine-precision limits, so certification is claimed only within the conditioning range the adversarial suite (well-conditioned, largely diagonal or near-diagonal mass matrices) actually exercises. The governing principle is stated plainly: conditioning can establish that an equality is unresolved; it can never establish that an equality holds. A cluster that shows no deficiency and whose members are not resolvably separated is reported as near-coalescent-unresolved rather than as semisimple.

The certified multiplicity domain is stated and enforced rather than assumed, and it is narrower than two earlier releases of this engine claimed. ChemSA certifies only simple roots and isolated algebraic-multiplicity-two clusters, decided from finite-precision evidence alone: order 2 is provably decidable this way, since two distinct eigenvalues never produce a nonzero rank-one nullity at their mean. NO other verdict is ever certified from floating-point evidence alone. A prior release carved out an exception for members equal to within the linear roundoff radius -- an apparently exact higher-multiplicity coincidence, such as a critically damped twin whose four poles all equal \(-1\) to machine precision -- reasoning that no alternative diagonalizable system could produce that pattern at every achievable precision. That reasoning was WRONG: a triangular matrix with three eigenvalues differing by \(10^{-15}\) and \(2\times10^{-15}\) is genuinely, exactly distinct and diagonalizable, since \(K=1-10^{-15}\) is an ordinary, exactly representable float64 stiffness rather than roundoff noise, yet it fell within that same radius and was certified defective with full confidence. Floating-point proximity, no matter how tight, is never proof of exact mathematical coincidence, so that carve-out has been removed. For a candidate of size \(m>2\), two independent tests are run: a Jordan-chain continuation test (confirming a generalized eigenspace of dimension exactly \(m\) with genuine deficiency) and a Puiseux-fan geometry test (confirming its members are spread the way \(m\) genuine roots of a common defect are). Passing both is NECESSARY evidence for a genuine order-\(m\) defect; it is NOT SUFFICIENT, and is reported only as \texttt{higher-order-ep-candidate}, never certified, REGARDLESS of how small the members' spread is -- a critically damped twin now reports as a candidate like everything else of its size. The reason is structural, not a numerical shortfall: a real, non-roundoff perturbation of a Jordan block, \(J_m(-1)\) with its \((m,1)\) entry offset by an actual \(\varepsilon\), has characteristic polynomial \((\lambda+1)^m=\varepsilon\) and therefore \(m\) genuinely distinct, fully resolvable, diagonalizable roots for every \(\varepsilon\neq0\) -- separated by as much as \(10^{-2}\) at \(m=6\) -- and both tests pass this system identically to a genuine defect, because the Puiseux fan they detect is produced by PERTURBING a Jordan block, and observing that fan does not prove the perturbed matrix remains one. No finite-precision test on a single matrix instance can settle that in general; only a structural or symbolic argument, or parameter continuation establishing that geometric multiplicity genuinely collapses along a swept control parameter, can. A candidate component is refused outright, and returned near-coalescent-unresolved without any label, whenever no plausible local partition can be established at all: a repeated pair not isolated from the nearest external pole by more than its clustering radius, or a component of more than two members that fails either validation test. This boundary is exercised by mandatory crowded-spectrum and order-aware regression suites (\texttt{test\_crowded\_spectrum.py}, \texttt{test\_order\_aware\_ep.py}) wired into the master validator, the latter verified against zero false certifications across perturbation scales from \(10^{-2}\) to \(10^{-12}\), including the triangular-matrix and distinct-scalar-QEP counterexamples that motivated this correction.

\subsection{Biorthogonal response layer}

Because eigenvalue-based readings cannot determine observables (two systems can share identical poles, \(\chi\), and pole motion yet produce different measured spectra), a response layer built from left and right eigenvectors supplies channel-specific residues, structured perturbation sensitivities, and declared-weight condition numbers. For isolated or uniquely partitioned nondefective clusters passing the stated separation and conditioning gates, the response layer reconstructs the supplied-coordinate resolvent, and refuses at defective poles, where the simple-pole residue is invalid rather than merely inaccurate. Crowded neighborhoods whose repeated subset cannot be uniquely isolated are returned near-coalescent-unresolved and excluded from response reconstruction rather than assigned a multiplicity: a candidate component with more than one plausible local partition (for example an exceptional point beside a nearby simple pole) is refused by an isolation and partition-uniqueness gate before any structure or response is reported. This boundary is enforced in code and exercised by mandatory crowded-spectrum regression tests (\texttt{test\_crowded\_spectrum.py}, wired into the master validator). The layer enforces the reciprocal mechanical input contract (\(M\) Hermitian positive definite, \(C\) and \(K\) Hermitian), so the two layers cannot accept different physical domains. Passivity is not assumed: negative damping and negative stiffness remain classifiable as active or unstable. The response layer accepts Hermitian complex coefficients; the classifier path is exercised and tested on real symmetric coefficients, and complex Hermitian input to the classifier is outside the tested domain. A second, physically observable signature of defectiveness follows from this layer: under a generic perturbation of size \(\varepsilon\), eigenvalues at a semisimple degeneracy split linearly in \(\varepsilon\), while at an EP2 (the only order this signature is checked against) they split as \(\sqrt{\varepsilon}\). Computed log-log slopes are 1.00 (semisimple) and 0.50 (defective), corroborating the multiplicity test of Section~2.3(a) through a completely different mechanism.

\subsection{External cross-check: an experimentally realized exceptional-point model class}

The demonstrations in Section~2.3 use constructed matrices, which is the right way to show what eigenvector deficiency means but does not test the engine against anything the world produced. That gap is closed with a system in which an exceptional point has been realized and measured in a laboratory: two linearized spring-coupled pendulums with two real controllable parameters, the length of one pendulum (moving the mass and stiffness matrices together) and a viscous-like damping from an eddy-current brake (moving the damping matrix alone) \citep{even2024}. That work reports, independently, coalescence of both eigenvalues and eigenvectors at the exceptional point, detection of eigenvector coalescence through the condition number, and Puiseux (fractional-exponent) eigenvalue behaviour in its neighbourhood. It also discriminates the point by comparing a pure-exponential fit of its measured response against a Jordan-chain fit, which is the same secular criterion applied in the time domain. Every constant used here is the authors' own calibrated value, identified by them from their experiment and taken at the full precision released in their software \citep{evencode}. The exceptional point is then located by minimizing the eigenvalue gap over the same two real controllable parameters. That search is a local refinement initialized in the neighbourhood of the published coordinates, not an answer-blind global search, and it is reported as such; what makes it informative is that the published six-figure coordinates do not themselves return a degenerate pair. It locates the point at pendulum length \(l_2 = 0.674843\)~m and brake damping \(c_2 = 0.132988\)~N\,m\,s (residual eigenvalue gap \(4.4\times10^{-9}\)), against their published \(0.674831\)~m and \(0.133028\)~N\,m\,s: agreement to \(1.2\times10^{-5}\)~m and \(4.0\times10^{-5}\)~N\,m\,s. The eigenvalue there, \(-0.023143 - 3.271192i\), matches their reported \(-0.023155 \pm 3.271194i\) to \(2.5\times10^{-6}\)~rad\,s\(^{-1}\) in frequency. Evaluating instead at their published six-figure coordinates returns two simple poles separated by about \(10^{-3}\), which is not a disagreement but the codimension-two Puiseux sensitivity the source itself emphasizes: six figures do not place a system on an exceptional point, so it must be relocated in the same parameters. The reported signatures reproduce: the structure verdict is defective-at-tolerance; algebraic multiplicity 2 exceeds geometric multiplicity 1; the declared-weight condition number diverges on approach and is refused at the exceptional point itself. The pendulum validates the mathematics (a quadratic eigenvalue problem solved by companion linearization); a phonon application additionally requires the reciprocal mechanical input contract of Section~2.5, which Hermitian force-constant and damping matrices satisfy. The authors' measured time traces are not ingested, so this reproduces their calibrated model rather than their raw data, and it is not an independent experimental validation. What it establishes is narrower and sufficient: that the machinery fires on a physically parameterized, laboratory-realizable model class located by search, rather than only on a constructed Jordan block.
